\section{A deterministic rank bound}\label{rank:section}

We now turn expansion of a finite system of relators into a lower bound
for the number of generators of its presented group.  The proof represents
a generating family by a labelled graph.  After absorbing a small set of
letters into a coefficient group, Lemma~\ref{planar:main} finds a long
relator segment in that graph.  Replacing this segment by the short
complement of the relator contradicts minimality of the graph.
This is an Arzhantseva--Ol'shanskii shortening argument
\cite[\S~Marked graphs]{arzhantseva-olshanskii-1996}; see also
\cite[Definitions~2.6--2.9 and Proposition~2.11]{kapovich-schupp-2002}.
The coefficient group and the saturation procedure below allow us to
use the move for the present system of relators.

\begin{theorem}\label{rank:main}
Let $V$ and $\mathcal T$ be sets of cardinality $n$.  For each
$t\in\mathcal T$, let
\[
 r_t=x_{c(t,0)}x_{c(t,1)}\cdots x_{c(t,99)},\qquad c(t,i)\in V,
 \qquad
 D=\langle x_v\ (v\in V)\mid r_t=1\ (t\in\mathcal T)\rangle.
\]
For $T\subseteq\mathcal T$, put
$N(T)=\{c(t,i):t\in T,\ 0\le i\le99\}$.
Call a row $t$ bad if it has a repeated column, or if it shares at least
two distinct columns with another row, and let $b$ be the number of bad
rows.  Suppose that
\[
 0<\alpha<\frac1{200},\qquad \alpha n\ge1,
 \qquad 100b\le\alpha n,
\]
and that
\begin{equation}\label{rank:expansion}
 |N(T)|>96|T|\qquad
 \text{whenever }1\le |T|\le\alpha n.
\end{equation}
Then the minimum size of a generating set of $D$ satisfies
\[
 \operatorname{rk}(D)>\frac{\alpha n}{100}.
\]
\end{theorem}

\begin{proof}
Suppose instead that $D$ has at most $\alpha n/100$ generators.
We proceed in four steps: choose a minimal graph, saturate a small
coefficient alphabet, locate a long relator segment, and shorten the graph.

\subsection*{A minimal labelled graph}
Give each edge of a finite graph an orientation and a label $x_v$ or
$x_v^{-1}$; reading the edge backwards inverts its label.  Path labels
are words in these letters, evaluated in $D$.  A generating set written
as words in the $x_v$ gives a subdivided rose whose based loops map onto
$D$.  Among all finite connected labelled graphs $Y$ with
\[
 \beta(Y)=|E(Y)|-|V(Y)|+1\le\frac{\alpha n}{100}
\]
and surjective loop-label map to $D$, choose one with the fewest edges.
Write $\beta=\beta(Y)$.  Changing basepoint conjugates the loop image in
$D$, so it preserves surjectivity.  Consequently $Y$ has no vertex of
degree one: pruning such a vertex and its edge, and moving the basepoint
if necessary, would give a smaller graph with the same property.

Choose a basepoint $p$.  Decompose $Y$ into arcs by splitting at $p$ and
at every vertex of degree at least three.  An arc is an edge path with
distinct interior vertices, all of degree two and different from $p$;
its endpoints may coincide.  If $Y$ is a point, it has no arcs.  If $Y$
is a circle, then $\beta=1$ and marking $p$ gives one arc.  In all other
cases, the identity
\[
 \sum_{v\in V(Y)}(\deg(v)-2)=2\beta-2
\]
bounds the number of vertices of degree at least three by $2\beta-2$.
Suppressing all degree-two vertices except possibly $p$ leaves at most
$2\beta-1$ vertices and therefore at most $3\beta-2$ edges.
In particular, in every case the number of arcs is at most $3\beta$.

\subsection*{Saturating the coefficient alphabet}
Start with $S\subseteq V$ equal to the set of columns occurring in bad
rows.  Apply either of the following rules whenever it adds a new column:
\begin{enumerate}
\item If an arc has at most $30$ occurrences of labels with index
outside $S$, put every index occurring on that arc into $S$.
\item If a row has at least $10$ occurrences with index in $S$, put
all its columns into $S$.
\end{enumerate}
The process terminates, since each step enlarges $S\subseteq V$.
An arc can contribute at most $30$ new columns and is used at most once.
Thus the initial set and all arc additions together contribute at most
\begin{equation}\label{rank:seed}
 100b+30(3\beta)
 \le\alpha n+\frac{90\alpha n}{100}<2\alpha n.
\end{equation}
Every nontrivial row addition contributes at most $90$ new columns.

Put $s=\lfloor\alpha n\rfloor\ge1$.  If $s$ rows were processed by
the second rule, let $T$ consist of the first $s$ such rows.  At that
moment,
\[
 N(T)\subseteq S,
 \qquad |S|\le2\alpha n+90s<96s.
\]
For the strict inequality, use
$\alpha n<s+1\le2s<3s$.  This contradicts
\eqref{rank:expansion}.  Thus fewer than $s$ rows are processed, and the
final set satisfies
\begin{equation}\label{rank:saturation}
 |S|\le2\alpha n+90(s-1)\le92\alpha n<n.
\end{equation}
Let $U=V\setminus S$, which is nonempty.  At termination, every arc
has either no $U$-labelled edge or more than $30$ such occurrences.
Moreover, every row not entirely in $S$ is good and has at most nine
$S$-occurrences.  Such a row has at least $91$ distinct positive
$U$-letters, and distinct such rows share at most one $U$-letter.

Let $H=\langle x_s:s\in S\rangle$ be the actual subgroup of $D$
generated by these elements.  There is an isomorphism
\begin{equation}\label{rank:relative-presentation}
 D\cong
 \bigl(H*F(x_u:u\in U)\bigr)
 \big/\langle\!\langle r_t:\text{$r_t$ is not entirely in $S$}
 \rangle\!\rangle.
\end{equation}
Indeed, the group on the right maps to $D$ because every defining
relation of $H$ and every displayed row hold in $D$.  Conversely, its
generators satisfy all the original rows: an omitted row is a true
relation of $H$, and every other row is retained.  The resulting maps
are inverse on the generators, hence are inverse isomorphisms.
In particular, the coefficient group in
\eqref{rank:relative-presentation} embeds in the quotient.  No finite
presentation of $H$ is being assumed.

\subsection*{Finding a long segment in a shortest loop}
Call an occurrence of $x_u^{\pm1}$ with $u\in U$ exterior.
Choose $j\in U$ and a shortest edge loop $P$ based at $p$ whose label
represents $x_j$ in $D$.  This includes the possibility that $P$ is
empty.  It has no immediate backtracks and therefore traverses complete
arcs.  Its number of exterior occurrences is consequently either zero
or greater than $30$.

The label of $P$ is linearly $H$-reduced: two consecutive exterior
occurrences in its linear list cannot be inverse letters separated by
an $S$-word trivial in $H$.  To prove this, suppose that the relevant
portion of $P$ is $eRf$, where $R$ uses only $S$-edges and has value
$1$ in $H$.  If $f=e^{-1}$ as oriented edges, the closed detour $eRf$
has value $1$ in $D$, so deleting it shortens $P$ without changing its
value.  Otherwise $e$ and $f$ are different unoriented edges, and $R$
uses neither of them.  Slide the initial incidence of $f$ backwards
along $R$ to the terminal vertex of $e$.  This keeps the numbers of
vertices and edges unchanged and keeps the graph connected: the path
$R$ is still present.  Each old traversal of $f$ can be replaced by
$R^{-1}f$ in the slid graph, and a reverse traversal by its inverse.
Since $R$ has value $1$ in $D$, these substitutions preserve the values
of all based loops, so the loop-label image still surjects onto $D$.
Now $e^{-1}$ and $f$ leave the same vertex with the same label.  Fold
them together.  This removes one edge and identifies either zero or
one pair of distinct vertices; the graph remains connected, its rank
does not increase, and the quotient preserves all loop labels.
It is therefore a smaller admissible graph, contrary to the choice of
$Y$.  This proves linear $H$-reduction.

Put $L=H*F(x_u:u\in U)$.  The normal form in this free product shows
that a linearly $H$-reduced word representing $x_j$ in $L$ has exactly
one exterior occurrence: after evaluating its $H$-gaps, adjacent free
letters separated by trivial gaps form freely reduced blocks.
Since $P$ has zero or more than $30$ exterior occurrences, its label
does not equal $x_j$ in $L$.
Thus $\operatorname{lab}(P)x_j^{-1}$ is nontrivial in $L$ but trivial
in $D$.

Cyclically $H$-reduce this word by cancelling adjacent inverse exterior
letters across gaps trivial in $H$.  Let $W$ be the resulting cyclic
word.  It still has an exterior occurrence.  Otherwise the original
word would be conjugate in $L$ to an element of $H$ trivial in $D$;
the embedding of $H$ would make that element, and hence the original
word, trivial in $L$.

There is just one possible disruption of the original path when
passing to $W$.  If no cancellation occurs, it is the appended letter
$x_j^{-1}$, which we call the seam letter.  If cancellation occurs,
the first cancellation must involve that appended letter, by linear
$H$-reduction of $P$.  All later cancellations occur between the two
ends of the remaining exterior list of $P$.  The surviving exterior
occurrences form a contiguous block of that list, all its intervening
gaps are unchanged, and only its closing gap, called the seam gap,
may have changed.

Apply Lemma~\ref{planar:main} to
\eqref{rank:relative-presentation} and $W$.  Each mixed row has at
least $91$ exterior occurrences.  If the lemma gives two
exterior-disjoint intervals, at least one avoids the seam letter,
when present.  When there is a seam gap, at most one interval can
cross it, because a crossing uses both adjacent exterior occurrences.
In either case an interval with at least $91-3=88$ exterior
occurrences remains intact in $P$.  If the lemma instead gives an
entire cyclic relator, remove the seam letter, when present, or cut at
the seam gap.  This leaves at least $90$ exterior occurrences intact
in $P$.

We have therefore obtained a subpath $Q$ of $P$, beginning and ending
with exterior edges, that matches a consecutive portion of a cyclic
row or its inverse and has at least $88$ exterior occurrences.
Its exterior indices are distinct.  Its intervening $S$-words have
the same values in $H$ as the corresponding row segments, although
their lengths need not agree.  The matching row segment has at least
$88$ of the original $100$ letters.  Its complementary segment has
length at most $12$, so the inverse complementary word gives a word
$z$ in the original generators such that
\begin{equation}\label{rank:short-word}
 [z]_D=[\operatorname{lab}(Q)]_D,
 \qquad 1\le |z|\le12.
\end{equation}
If the complementary word is empty, use $x_jx_j^{-1}$ for $z$.

\subsection*{Shortening the graph}
Decompose $Q$ according to its arc traversals, allowing partial
traversals at its two ends.  Some such traversal contains more than
$30$ exterior edges of $Q$.  Otherwise no complete traversal could
have exterior edges, by saturation, and the two partial traversals
would contribute at most $60$ exterior occurrences in total.
In the chosen traversal, let $I$ run from its first exterior edge
through its last exterior edge, inclusive.  Then
\[
 Q=Q_1 I Q_2,\qquad \ell:=|E(I)|\ge31.
\]
The segment $I$ has distinct interior vertices of degree two, none
equal to $p$; its endpoints may coincide.

The paths $Q_1,Q_2$ avoid both the edges and the interior vertices of
$I$.  First, neither endpoint of $Q$ is interior to $I$: at such a
vertex the first or last edge of $Q$ would be one of the edges of $I$,
because that vertex has degree two.  That edge is exterior, so its
two occurrences in $Q$ would repeat an exterior index.  Next, any
path entering the interior of $I$ from outside must use its first
or last edge.  Both are exterior and already occur in $I$, so neither
$Q_1$ nor $Q_2$ can enter this way.  These arguments also exclude a
second traversal of an edge of $I$; when $I$ is closed, the same two
boundary incidences give the identical conclusion.

Figure~\ref{rank:shortening-schematic} illustrates the construction.

\begin{figure}[ht]
\centering
\begin{tikzpicture}[>=Stealth, line width=.7pt,
  every node/.style={font=\small},
  endpoint/.style={circle,fill=black,inner sep=1.5pt},
  pathlabel/.style={fill=white,inner sep=2pt}]
  \begin{scope}
    \node at (2.45,2.15) {Before};
    \coordinate (a) at (0,0);
    \coordinate (b) at (1.35,0);
    \coordinate (c) at (3.55,0);
    \coordinate (d) at (4.90,0);
    \draw[->] (a) -- node[pathlabel,above] {$Q_1$} (b);
    \draw[->,very thick] (b) -- node[pathlabel,above] {$I$} (c);
    \draw[->] (c) -- node[pathlabel,above] {$Q_2$} (d);
    \foreach \q in {a,b,c,d} \node[endpoint] at (\q) {};
    \node[below=5pt] at (a) {$q_-$};
    \node[below=5pt] at (b) {$i_-$};
    \node[below=5pt] at (c) {$i_+$};
    \node[below=5pt] at (d) {$q_+$};
    \node at (2.45,-.88) {$Q=Q_1IQ_2,\qquad |E(I)|\ge31$};
  \end{scope}
  \begin{scope}[xshift=6.2cm]
    \node at (2.45,2.15) {After};
    \coordinate (a) at (0,0);
    \coordinate (b) at (1.35,0);
    \coordinate (c) at (3.55,0);
    \coordinate (d) at (4.90,0);
    \draw[->] (a) -- node[pathlabel,above] {$Q_1$} (b);
    \draw[->] (c) -- node[pathlabel,above] {$Q_2$} (d);
    \draw[->,very thick] (a) to[out=65,in=115]
      node[pathlabel,above] {$Z,\quad |E(Z)|\le12$} (d);
    \foreach \q in {a,b,c,d} \node[endpoint] at (\q) {};
    \node[below=5pt] at (a) {$q_-$};
    \node[below=5pt] at (b) {$i_-$};
    \node[below=5pt] at (c) {$i_+$};
    \node[below=5pt] at (d) {$q_+$};
    \node at (2.45,-.88) {Bypass from $i_-$ to $i_+$: $Q_1^{-1}ZQ_2^{-1}$};
  \end{scope}
\end{tikzpicture}
\caption{The shortening move: add $Z$ between the endpoints of $Q$,
then delete the edges and interior vertices of $I$. The retained paths
$Q_1,Q_2$ and the new path give a bypass with the same value in $D$
as $I$. The drawing is schematic: displayed endpoints may coincide,
and $Q_1,Q_2$ may overlap or be empty. The proof establishes that both
avoid the edges and interior vertices of $I$. Each line denotes a path,
and the rest of the graph is omitted.}
\label{rank:shortening-schematic}
\end{figure}


Attach a new path $Z$ between the endpoints of $Q$ with label $z$,
using new interior vertices and $k=|z|\le12$ new edges.  Then delete
the edges and interior vertices of $I$.  The old endpoints of $I$
are still connected by the path
\[
 Q_1^{-1} Z Q_2^{-1},
\]
which avoids everything deleted.  Thus the new graph is connected.
The inserted path has $k-1$ new interior vertices, and the deleted
segment has $\ell-1$ interior vertices, also when either path has
coincident endpoints.  Hence the changes in the numbers of edges
and vertices are both $k-\ell$, and the graph rank is unchanged.
The basepoint survives, since it is not an interior vertex of $I$.

Equation~\eqref{rank:short-word} gives
\[
 [\operatorname{lab}(Q_1^{-1}ZQ_2^{-1})]_D
   =[\operatorname{lab}(I)]_D.
\]
Every reduced loop at $p$ in the old graph that uses an edge of $I$
must traverse all of $I$, in one direction or the other: its interior
vertices have degree two and do not contain $p$.  Replacing each
such traversal by the displayed bypass, or its inverse, gives a loop
in the new graph with the same value in $D$.  The loop-label map of
the new graph therefore remains surjective.  Its edge count has
decreased by at least $31-12=19$, contradicting minimality of $Y$.
This proves the theorem.
\end{proof}

\begin{corollary}\label{rank:asymptotic}
Suppose that a sequence of length-$100$ row systems satisfies
\eqref{rank:expansion} for a fixed $0<\alpha<1/200$, their sizes
$n$ tend to infinity, and their bad-row counts satisfy $b=o(n)$.
Then their presented groups satisfy
\[
 \operatorname{rk}(D)>\frac{\alpha n}{100}
 \quad\text{for all sufficiently large }n,
 \qquad
 \liminf_{n\to\infty}\frac{\operatorname{rk}(D)}n
 \ge\frac\alpha{100}.
\]
\end{corollary}
\begin{proof}
Eventually $\alpha n\ge1$ and $100b\le\alpha n$, so
Theorem~\ref{rank:main} applies.
\end{proof}
